% SolarMax — Install Git (Homebrew) and use it inside VS Code (macOS)
% Compile:  pdflatex git_setup_mac.tex   (run twice for the table of contents)

\documentclass[11pt]{article}
\usepackage[margin=2.2cm]{geometry}
\usepackage{xcolor}
\usepackage[most]{tcolorbox}
\usepackage{enumitem}
\usepackage{titlesec}
\usepackage{listings}
\usepackage[colorlinks=true,linkcolor=blue!60!black,urlcolor=blue!60!black]{hyperref}

% ── Colors / section styling ─────────────────────────────────────────────────
\definecolor{brandblue}{HTML}{1565C0}
\definecolor{brandgreen}{HTML}{2E7D32}
\titleformat{\section}{\Large\bfseries\color{brandblue}}{\thesection.}{0.5em}{}

% ── Callout boxes ─────────────────────────────────────────────────────────────
\newtcolorbox{tip}{colback=green!6,colframe=brandgreen,boxrule=0.8pt,
  left=8pt,right=8pt,top=6pt,bottom=6pt,title=\textbf{Tip},fonttitle=\color{white}}
\newtcolorbox{warn}{colback=red!5,colframe=red!70!black,boxrule=0.8pt,
  left=8pt,right=8pt,top=6pt,bottom=6pt,title=\textbf{Important},fonttitle=\color{white}}
\newtcolorbox{note}{colback=blue!5,colframe=brandblue,boxrule=0.8pt,
  left=8pt,right=8pt,top=6pt,bottom=6pt,title=\textbf{Note},fonttitle=\color{white}}

% ── Terminal command blocks ──────────────────────────────────────────────────
\lstdefinestyle{cmd}{basicstyle=\ttfamily\small,backgroundcolor=\color{black!6},
  frame=single,rulecolor=\color{black!25},framesep=6pt,aboveskip=6pt,belowskip=6pt,
  xleftmargin=4pt,xrightmargin=4pt,breaklines=true,columns=fullflexible,keepspaces=true}
\lstset{style=cmd}

\setlist[enumerate]{leftmargin=1.7em,itemsep=3pt,topsep=4pt}

\begin{document}

% ── Title ─────────────────────────────────────────────────────────────────────
\begin{center}
{\Huge\bfseries\color{brandblue} SolarMax}\\[4pt]
{\LARGE Install Git \& Get the Code in VS Code --- macOS}\\[8pt]
{\large Install Homebrew \(\rightarrow\) install Git \(\rightarrow\) clone the project, all from VS Code.}\\[6pt]
\rule{0.9\linewidth}{0.8pt}
\end{center}

\vspace{4pt}
\tableofcontents
\vspace{8pt}

% =============================================================================
\section{What this does}
This installs \textbf{Git} on your Mac using \textbf{Homebrew} (a package installer),
adds a Git extension to \textbf{VS Code} so it's easy to see, and clones the SolarMax
project straight from inside VS Code. After this you get the latest code with one
click any time --- no re-downloading ZIPs.

\begin{note}
You already have VS Code + PlatformIO from the Flashing guide. We'll do everything
here \textbf{inside VS Code}, including the Terminal commands.
\end{note}

% =============================================================================
\section{What you need}
\begin{itemize}[leftmargin=1.6em]
  \item A Mac with internet, and VS Code installed.
  \item Your Mac \textbf{login password} (the installers ask for it).
\end{itemize}

% =============================================================================
\section{Step 1 --- Install Homebrew (in VS Code's terminal)}
\begin{enumerate}
  \item Open VS Code. In the top menu: \textbf{Terminal \(\rightarrow\) New Terminal}
        (a command panel opens at the bottom).
  \item Go to \href{https://brew.sh}{brew.sh} and copy its install command --- or just
        paste this one (it's the same) and press \texttt{Return}:
\begin{lstlisting}
/bin/bash -c "$(curl -fsSL https://raw.githubusercontent.com/Homebrew/install/HEAD/install.sh)"
\end{lstlisting}
  \item Type your \textbf{Mac password} when asked (the characters stay hidden --- that's
        normal) and press \texttt{Return}. Let it run; it takes a few minutes.
\end{enumerate}

\begin{note}
Homebrew may first install Apple's ``Command Line Tools'' automatically --- that's
expected, just let it finish.
\end{note}

\begin{warn}
When it's done, Homebrew prints a \textbf{``Next steps''} section with \textbf{two
lines} to run (they add Homebrew to your PATH). \textbf{Copy those two lines from your
terminal and run them.} On Apple-silicon Macs they look like this:
\begin{lstlisting}
echo 'eval "$(/opt/homebrew/bin/brew shellenv)"' >> ~/.zprofile
eval "$(/opt/homebrew/bin/brew shellenv)"
\end{lstlisting}
\end{warn}

Confirm Homebrew works:
\begin{lstlisting}
brew --version
\end{lstlisting}
A version number (e.g.\ \texttt{Homebrew 4.x.x}) means it's ready.

% =============================================================================
\section{Step 2 --- Install Git with Homebrew}
In the same terminal:
\begin{lstlisting}
brew install git
\end{lstlisting}
Then confirm:
\begin{lstlisting}
git --version
\end{lstlisting}
Any version number (e.g.\ \texttt{git version 2.45.0}) means Git is installed.

\begin{tip}
Now \textbf{quit and reopen VS Code} once. This lets it detect the newly installed
Git and turn on its Source Control features.
\end{tip}

% =============================================================================
\section{Step 3 --- Add the GitLens extension in VS Code}
VS Code already has built-in Git support; \textbf{GitLens} makes it more visual (history,
who changed what, a commit graph) --- nice for beginners.
\begin{enumerate}
  \item Click the \textbf{Extensions} icon in the left bar (four squares).
  \item Search for \texttt{GitLens}.
  \item Click \textbf{Install} on ``GitLens --- Git supercharged''.
\end{enumerate}

\begin{note}
This is optional but recommended. Even without it, the \textbf{Source Control} icon in
the left bar (the branch symbol) already lets you clone, pull, and commit.
\end{note}

% =============================================================================
\section{Step 4 --- Clone SolarMax from inside VS Code}
No terminal needed for this --- use VS Code's command palette:
\begin{enumerate}
  \item Press \texttt{Cmd}\,+\,\texttt{Shift}\,+\,\texttt{P} to open the Command Palette.
  \item Type \texttt{Git: Clone} and select it.
  \item Paste the project's address and press \texttt{Return}:
\begin{lstlisting}
https://github.com/zhsquared/SolarMax.git
\end{lstlisting}
  \item Choose a folder to put it in (your \textbf{Documents} folder is a good choice)
        and click \textbf{Select as Repository Destination}.
  \item When VS Code asks \textbf{``Would you like to open the cloned repository?''},
        click \textbf{Open}.
\end{enumerate}
You now have the newest SolarMax code open in VS Code.

\begin{warn}
If you had an older \textbf{ZIP} copy of the project, use this cloned folder from now
on, and re-apply your own \texttt{config.h} settings here --- your Wi-Fi (lines 9--10)
and location (lines 4--5).
\end{warn}

% =============================================================================
\section{Step 5 --- Getting updates later (in VS Code)}
When there's new code, pull it without re-downloading anything:
\begin{enumerate}
  \item Click the \textbf{Source Control} icon in the left bar (branch symbol).
  \item Click the \textbf{``\ldots''} menu at the top of that panel
        \(\rightarrow\) \textbf{Pull}.
\end{enumerate}
(Or click the little \textbf{sync} arrows next to the branch name at the bottom-left of
the window.) That's it --- you're on the newest version.

% =============================================================================
\section{Troubleshooting}
\begin{itemize}[leftmargin=1.6em]
  \item \textbf{\texttt{brew: command not found}:} the PATH step didn't run. Re-run the
        two ``Next steps'' lines Homebrew printed (Step~1), then \textbf{reopen the
        terminal} and try \texttt{brew --version} again.
  \item \textbf{\texttt{git: command not found} after installing:} quit and reopen VS
        Code (and the terminal) so it picks up the new install.
  \item \textbf{Password prompt shows nothing as you type:} that's normal for terminal
        passwords --- type it and press \texttt{Return}.
  \item \textbf{\texttt{Git: Clone} isn't in the Command Palette:} Git isn't detected
        yet --- finish Step~2, then quit and reopen VS Code.
  \item \textbf{Clone fails partway:} usually a network hiccup --- run \texttt{Git:
        Clone} again with the same address.
\end{itemize}

\end{document}
